\section{Obyčejné diferenciální rovnice}
\subsection{Základní pojmy}

\begin{definice}
Diferenciální rovnicí rozumíme rovnici tvaru
\begin{equation}\label{odrdef}
F(x,y,y',y'',\ldots,y^{(n-1)},y^{(n)})=0.
\end{equation}
kde $F$ je reálná funkce $n+2$ proměnných.

\textbf{Řešením diferenciální rovnice} (\ref{odrdef}) rozumíme funkci $y$ definovanou na nějakém neprázdném
otevřeném intervalu $I$, která má v každém bodě intervalu $I$ vlastní $n$-tou derivaci a jejíž hodnoty
spolu s hodnotami derivací splňují rovnici (\ref{odrdef}) v každém bodě intervalu $I$, tj. pro každé $x\in I$
platí
\[
F\bigl(x,y(x),y'(x),y''(x),\ldots,y^{(n-1)}(x),y^{(n)}(x)\bigr)=0.
\]
\end{definice}

\begin{priklad}
	Rovnice $p^\prime = ap - bp^2$ nebo $y^{\prime\prime} +y^\prime +\ = \sin x$. Tu lze převést na
$$y^{\prime\prime} +y^\prime +\ -\sin x= 0.$$
Tady je
	$$F(z_1,z_2,z_3,z_4)=--\sin z_1 + z_2 + z_3 + z_4.$$
\end{priklad}

\begin{definice}
Je-li funkce $y$ řešením rovnice (\ref{odrdef}) na intervalu $I$ a funkce $\widetilde y$ řešením rovnice (\ref{odrdef}) na
intervalu $\widetilde I$, kde $I\subset\widetilde I$, $I\ne\widetilde I$, a $y(x)=\widetilde y(x)$ pro všechna $x\in I$,
pak říkáme, že řešení $\widetilde y$ je \textbf{prodloužením řešení $y$} na interval $\widetilde I$.
Řešení rovnice (\ref{odrdef}), které nemá prodloužení, nazýváme \textbf{maximálním řešením} rovnice (\ref{odrdef}).
\end{definice}

\subsection{Rovnice se separovanými proměnnými}
\begin{definice}
Rovnice se separovanými proměnnými je rovnice tvaru
\begin{equation}\label{separ}
y'=g(y)h(x).
\end{equation}
kde $g,h$ jsou dané reálné funkce.
\end{definice}

\begin{pozorovani}
Označme $\mathcal D(h)=I$.
	Předpokládejme, že $g(c)=0$. Položme $y(x)=c$, $x \in I$. To je maximální řešení, neboť
	po dosazení je na levé straně $y^\prime(x)=0$ a na pravé straně
	$$
	    g(y(x))\cdot h(x) = g(c)\cdot h(x) = 0\cdot h(x)=0.
	$$
\end{pozorovani}

\begin{pozorovani}
Předpokládejme, že $J$ je otevřený interval, kde $g$ nemění znaménko a $y: L \to J$ řeší na otevřeném
	intervalu $L\subseteq I$ rovnici (\ref{separ}). Potom
	\begin{align*}
	\forall x \in L: \qquad y^\prime (x) &= g(y(x))\cdot h(x),\\
	\frac{y^\prime(x)}{g(y(x))}&= h(x).
	\end{align*}
	Nechť $H$ je primitivní k funkci $h$, tj. $H^\prime(x) = h(x)$, $x \in I$. Nechť $G$ je primitivní
	k funkci $\frac{1}{g}$, tj. $G^\prime(z)=\frac{1}{g(z)}$, $z\in J$.
	Protože je
	\begin{align*}
	(G(y(x)))^\prime &= \frac{1}{g(y(x))}\cdot y^\prime(x), \qquad x \in L \\
	H^\prime(x)&=h(x), \qquad x \in L.
	\end{align*}
	Existuje tedy $k \in \mathbb R $ takové, že
	$$
	    \forall x \in L: \qquad G(y(x)) = H(x) + k.
	$$
	Pak je
	$$
	    y(x) = G^{-1} (H(x)+k).
	$$
\end{pozorovani}

\paragraph{Metoda řešení}
Předpokládejme, že $\mathcal D(g)$ je interval, $\mathcal D(h)$ je otevřený interval, $g$ je spojitá na
$\mathcal D(g)$, a $h$ je spojitá
na $\mathcal D(h)$.

\paragraph{1. krok (kořeny a znaménko $g$).}
Označme $I=\mathcal D(h)$. Nalezneme všechny nulové body funkce $g$.
Je-li totiž $g(c)=0$, pak funkce $y(x)=c$, $x\in I$, je řešením rovnice (\ref{separ}).
Těmto řešením se říká \textbf{singulární nebo také stacionární}. Určíme maximální otevřené intervaly,
na kterých je $g$ nenulová.

\paragraph{2. krok (výpočet primitivních funkcí).}
Nechť $J$ je interval z 1. kroku, kde je $g$ nenulová.
Budeme hledat řešení, která jsou definovaná na nějakém otevřeném intervalu, který je obsažen
v $I$, a mají hodnoty v intervalu $J$. Nalezneme primitivní funkci $H$ k funkci $h$ na $I$ a primitivní
funkci $G$ k funkci $\frac{1}{g}$ na $J$. Existuje inverzní funkce
$G^{-1}\colon G(J)\to J$.

\paragraph{3. krok (určení definičních oborů).}
Zvolme $c\in\mathbb R$. Nalezneme maximální neprázdné otevřené intervaly obsažené v množině
\[
\{x\in I;\ H(x)+c\in G(J)\}.
\]

Nechť $L$ je takový interval. Na $L$ definujeme funkci $y$ předpisem
\[
y(x)=G^{-1}\bigl(H(x)+c\bigr),\qquad x\in L.
\]
Funkce $y$ je dobře definována a je řešením rovnice (\ref{separ}).

Ověření, že funkce $y(x)$ je řešením rovnice (\ref{separ}).

Na levé straně je
\begin{align*}
	y^\prime(x) &= \left( G^{-1}(H(x)+c) \right ) ^\prime = \frac{1}{G^\prime \left( G^{-1}(H(x)+c) \right ) }\cdot (H(x)+c)\\
	&=g ( \underbrace{G^{-1}(H(x)+c)}_{y(x)} )\cdot h(x) = g(y(x))\cdot h(x).
\end{align*}

\paragraph{4. krok (lepení).}
Z řešení nalezených ve 3. kroku a singulárních řešení z 1. kroku slepíme všechna maximální řešení.
Nechť $y_1$ a $y_2$ jsou řešení rovnice (\ref{separ}), první na intervalu $(a,b)$ a druhé na intervalu $(b,c)$,
přičemž $b\in I$. Předpokládejme, že
\[
\lim_{x\to b-}y_1(x)=\lim_{x\to b+}y_2(x)=\alpha\in \mathcal D(g).
\]
Pak funkce
\[
y(x)=
\begin{cases}
y_1(x), & x\in(a,b),\\
\alpha, & x=b,\\
y_2(x), & x\in(b,c),
\end{cases}
\]
je řešením rovnice (\ref{separ}) na intervalu $(a,c)$.


Kontrola, že po čtvrtém kroku je v bodě $b$ nová funkce opravdu řešení. Zřejmě je (\ref{separ}) splněna pro každé
$x \in (a,b) \cup (b,c)$. Zbývá ověřit (\ref{separ}) pro $x=b$. Funkce $y$ je spojitá zleva v bodě $b$. Je
\begin{align*}
	y_-^\prime(b) &= \lim_{x \to b-} y^\prime(x) = \lim_{x \to b-} y_1^\prime(x) = \lim_{x \to b-} g(y_1(x))\cdot h(x)\\
	&= g(\alpha)\cdot h(b)=g(y(b))\cdot h(b).
\end{align*}
Z pravé strany provedeme obdobně.
